% Part: computability
% Chapter: recursive-functions
% Section: sequences

\documentclass[../../../include/open-logic-section]{subfiles}

\begin{document}

\olfileid{cmp}{rec}{seq}
\olsection{অনুক্রম}

আদিম পুনরাবৃত্ত অপেক্ষকসমূহের সেটটি উল্লেখযোগ্যভাবে শক্তিশালী। তবে
\emph{অনুক্রম} নিয়ে কাজ করার উপযুক্ত উপায় তৈরি হলে আমরা আরও অনেক কিছু
করতে পারব। স্বাভাবিক সংখ্যার সসীম অনুক্রমগুলিকে নিচের উপায়ে স্বাভাবিক
সংখ্যা দিয়ে চিহ্নিত করব: অনুক্রম $\langle a_0, a_1, a_2, \dots, a_k
\rangle$ নিচের সংখ্যাটির সঙ্গে সংশ্লিষ্ট:
\[
p_0^{a_0+1} \cdot p_1^{a_1+1} \cdot p_2^{a_2+1} \cdot \dots \cdot
p_k^{a_k+1}.
\]
সূচকগুলির সঙ্গে এক যোগ করা হয়, যাতে যেমন $\langle 2, 7, 3\rangle$ এবং
$\langle 2, 7, 3, 0, 0 \rangle$ অনুক্রম দুটির সাংখ্যিক সংকেত আলাদা হওয়া
নিশ্চিত হয়। $0$ ও~$1$—উভয়কেই শূন্য অনুক্রমের সংকেত ধরা যায়; নির্দিষ্ট
করে বলার জন্য $\emptyseq$ দিয়ে~$0$ বোঝানো যাক।

অনুক্রমের এই সংকেতায়ন যে কার্যকর, তার কারণ তথাকথিত পাটিগণিতের মৌলিক
উপপাদ্য: প্রত্যেক স্বাভাবিক সংখ্যা $n \ge 2$-কে একটিমাত্র উপায়ে নিচের
রূপে লেখা যায়:
\[
n = p_0^{a_0} \cdot p_1^{a_1} \cdot \dots \cdot p_k^{a_k}
\]
যেখানে $a_k \ge 1$। এটি নিশ্চিত করে যে চিত্রণ $\tuple{}(a_0,
\dots, a_k) = \tuple{a_0, \dots, a_k}$ একৈক: ভিন্ন অনুক্রম ভিন্ন
সংখ্যায় চিত্রিত হয়; প্রত্যেক সংখ্যার সঙ্গে সর্বাধিক একটি অনুক্রম
সংশ্লিষ্ট।

এখন আমরা দেখাব যে অনুক্রমের দৈর্ঘ্য নির্ণয়, তার $i$-তম উপাদান
নির্ণয়, অনুক্রমের শেষে একটি উপাদান যোগ এবং দুটি অনুক্রমের সংযুক্তি—সবই
আদিম পুনরাবৃত্ত।

\begin{prop}
  অনুক্রম $s$-এর দৈর্ঘ্য ফেরত দেয় এমন অপেক্ষক $\len{s}$ আদিম পুনরাবৃত্ত।
\end{prop}

\begin{proof}
  $R(i, s)$ সম্বন্ধটিকে এভাবে সংজ্ঞায়িত করা যাক:
  \[
  R(i, s) \text{ যদি এবং কেবল যদি }
  p_i \mid s \land p_{i+1} \nmid s.
  \]
  $R$ যে আদিম পুনরাবৃত্ত, তা স্পষ্ট। যখনই $s$ কোনো অশূন্য অনুক্রমের
  সংকেত, অর্থাৎ
  \[
  s = p_0^{a_0+1} \cdot \dots \cdot p_{k}^{a_{k}+1},
  \]
  তখন $p_i$ যদি এমন বৃহত্তম মৌলিক সংখ্যা হয় যার জন্য $p_i \mid s$,
  অর্থাৎ $i = k$ হলে $R(i,s)$ সত্য হয়। তাই $p_i$ যদি~$s$-কে নিঃশেষে ভাগকারী
  বৃহত্তম মৌলিক সংখ্যা হয়, তবে এবং কেবল তখনই $s$-এর দৈর্ঘ্য $i+1$;
  অতএব আমরা লিখতে পারি
  \[
  \len{s} =
  \begin{cases}
    0 & \text{যদি $s = 0$ অথবা $s = 1$} \\
    1 + \bmin{i < s}{R(i, s)} & \text{অন্যথায়।}
  \end{cases}
  \]
  এখানে সীমাবদ্ধ ন্যূনীকরণ ব্যবহার করা যায়, কারণ $s$ কোনো অনুক্রমের
  সংকেত হলে মাত্র একটি $i$ সম্বন্ধ $R(i,s)$ পূরণ করে, এবং এমন $i$ থাকলে
  তা~$s$-এর চেয়ে ছোট।
\end{proof}

\begin{prop}
  অনুক্রম $s$-এর শেষে $a$ যোগ করার ফল ফেরত দেয় এমন অপেক্ষক
  $\fn{append}(s,a)$ আদিম পুনরাবৃত্ত।
\end{prop}

\begin{proof}
  $\fn{append}$-কে এভাবে সংজ্ঞায়িত করা যায়:
  \[
  \fn{append}(s,a) =
  \begin{cases}
    2^{a+1} & \text{যদি $s = 0$ অথবা $s = 1$} \\
    s \cdot p_{\len{s}}^{a+1} & \text{অন্যথায়।}
  \end{cases}
  \]
\end{proof}

\begin{prop}
  $s$-এর $i$-তম উপাদান (যেখানে প্রথম উপাদানকে $0$-তম বলা হয়) ফেরত
  দেয়, অথবা $i$ যদি $s$-এর দৈর্ঘ্যের সমান বা বড় হয় তবে $0$ ফেরত দেয়,
  এমন অপেক্ষক $\fn{element}(s,i)$ আদিম পুনরাবৃত্ত।
\end{prop}

\begin{proof}
লক্ষ করুন, $a$ হলো~$s$-এর $i$-তম উপাদান যদি এবং কেবল যদি
$p_i^{a+1}$, $s$-কে নিঃশেষে ভাগকারী~$p_i$-এর বৃহত্তম ঘাত হয়; অর্থাৎ
$p_i^{a+1} \mid s$ কিন্তু $p_i^{a+2} \nmid s$। তাই:
  \[
  \fn{element}(s,i) =
  \begin{cases}
    0 & \mbox{যদি $i \geq \len{s}$} \\
    \bmin{a < s}{(p_i^{a+2} \nmid s)} & \text{অন্যথায়।}
  \end{cases}
  \]
\end{proof}

উপরে সংজ্ঞায়িত অপেক্ষকগুলির বিধিসম্মত নামের বদলে আমরা আরও সংক্ষিপ্ত
সংকেতলিপি ব্যবহার করব। $\fn{element}(s,i)$-এর বদলে $(s)_i$ লিখব, এবং
$\tuple{s_0, \dots, s_k}$ দিয়ে নিচের রাশিটি সংক্ষেপে বোঝাব:
\[
\fn{append}(\fn{append}(\dots \fn{append}(\emptyseq,s_0)
\dots),s_k).
\]
লক্ষ করুন, $s$-এর দৈর্ঘ্য~$k$ হলে $s$-এর উপাদানগুলি হলো
$(s)_0$, \dots,~$(s)_{k-1}$।

\begin{prop}
দুটি অনুক্রমকে সংযুক্ত করে এমন অপেক্ষক $\fn{concat}(s,t)$ আদিম
পুনরাবৃত্ত।
\end{prop}

\begin{proof}
  আমরা এমন একটি অপেক্ষক $\fn{concat}$ চাই, যার ধর্ম
  \[
    \fn{concat}(\tuple{a_0, \dots, a_k}, \tuple{b_0, \dots, b_l}) =
    \tuple{a_0, \dots, a_k, b_0, \dots, b_l}.
  \]
  আমরা একটি “সহায়ক” অপেক্ষক $\fn{hconcat}(s,t,n)$ ব্যবহার করব, যা
  $t$-এর প্রথম $n$টি প্রতীককে~$s$-এর সঙ্গে সংযুক্ত করে। আদিম
  পুনরাবৃত্তির সাহায্যে এই অপেক্ষকটিকে নিচের মতো সংজ্ঞায়িত করা যায়:
  \begin{align*}
    \fn{hconcat}(s,t,0) & = s\\
    \fn{hconcat}(s,t,n+1) & = \fn{append}(\fn{hconcat}(s,t,n),(t)_n)
    \intertext{তখন $\fn{concat}$-কে এভাবে সংজ্ঞায়িত করা যায়:}
    \fn{concat}(s,t) & = \fn{hconcat}(s,t,\len{t}).
  \end{align*}
\end{proof}

$\fn{concat}(s,t)$-এর বদলে আমরা $s \concat t$ লিখব।

কোনো অনুক্রমের দৈর্ঘ্য ও তার বৃহত্তম উপাদানের সাহায্যে অনুক্রমটির
সাংখ্যিক সংকেতকে সীমাবদ্ধ করতে পারা আমাদের কাজে লাগবে। ধরা যাক, $s$
দৈর্ঘ্য~$k$-এর একটি অনুক্রম, যার প্রত্যেক উপাদান কোনো সংখ্যা~$x$-এর
সমান বা ছোট। তাহলে~$s$-এর সর্বাধিক $k$টি মৌলিক উৎপাদক আছে, প্রত্যেকটি
সর্বাধিক~$p_{k-1}$, এবং~$s$-এর মৌলিক উৎপাদক বিশ্লেষণে প্রত্যেকটির ঘাত
সর্বাধিক $x+1$। অন্যভাবে বললে, আমরা যদি
\[
\fn{sequenceBound}(x,k) = p_{k-1}^{k \cdot (x+1)},
\]
সংজ্ঞায়িত করি, তবে উপরে বর্ণিত অনুক্রম~$s$-এর সাংখ্যিক সংকেত সর্বাধিক
~$\fn{sequenceBound}(x,k)$।

অনুক্রমের উপর এমন একটি সীমা থাকলে সীমাবদ্ধ অনুসন্ধানের সাহায্যে নতুন
অপেক্ষক সংজ্ঞায়িত করা যায়। যেমন, সীমাবদ্ধ অনুসন্ধানের সাহায্যে
$\fn{concat}$-কে সংজ্ঞায়িত করা যায়। আমরা যে বস্তুটি (সংযুক্ত অনুক্রমের
সংকেতসংখ্যা) খুঁজছি, তার একটি আদিম পুনরাবৃত্ত \emph{নির্দেশন} এবং কতদূর
খুঁজতে হবে তার একটি সীমা শুধু লিখে দিতে হবে। নিচেরটি কাজ করে:
\begin{align*}
  \fn{concat}(s,t) = {} & \bmin{v < \fn{sequenceBound}(s+t,\len{s} +
    \len{t})}{} \\
  & \quad(\len{v} = \len{s} + \len{t} \land {}\\
    & \qquad \bforall{i < \len{s}}{((v)_i = (s)_i) \land {} \\
      & \qquad \bforall{j < \len{t}}{((v)_{\len{s}+j} = (t)_j)})}
\end{align*}

\begin{prob}
দেখান যে এমন একটি আদিম পুনরাবৃত্ত অপেক্ষক~$\fn{sconcat}(s)$ আছে, যার
ধর্ম
\[
\fn{sconcat}(\tuple{s_0, \dots, s_k}) = s_0 \concat \dots \concat s_k.
\]
\end{prob}

\begin{prob}
দেখান যে এমন একটি আদিম পুনরাবৃত্ত অপেক্ষক~$\fn{tail}(s)$ আছে, যার ধর্ম
\begin{align*}
  \fn{tail}(\emptyseq) & = 0 \text{ এবং}\\
  \fn{tail}(\tuple{s_0, \dots, s_{k}}) & = \tuple{s_1, \dots, s_{k}}.
\end{align*}
\end{prob}

\begin{prop}
  \ollabel{prop:subseq}
  অপেক্ষক $\fn{subseq}(s, i, n)$, $s$-এর $i$-তম উপাদান থেকে শুরু হওয়া
  দৈর্ঘ্য~$n$-এর উপ-অনুক্রমটি ফেরত দেয়; এই অপেক্ষক আদিম পুনরাবৃত্ত।
\end{prop}

\begin{proof}
  অনুশীলনী।
\end{proof}

\begin{prob}
  \olref[cmp][rec][seq]{prop:subseq} প্রমাণ করুন।
\end{prob}

\end{document}
